\pdfinfoomitdate=1
\pdftrailerid{}
\pdfsuppressptexinfo=15
\documentclass[11pt,letterpaper]{article}
\usepackage[margin=1in]{geometry}
\usepackage[T1]{fontenc}
\usepackage{lmodern,amsmath,amssymb,amsthm,mathtools,booktabs,microtype}
\usepackage[colorlinks=true,linkcolor=black,citecolor=black,urlcolor=blue]{hyperref}
\usepackage{enumitem}
\setlist{itemsep=3pt,topsep=6pt}
\setlength{\parskip}{5pt}\setlength{\parindent}{0pt}
\setlength{\emergencystretch}{2em}
\newtheorem{theorem}{Theorem}[section]
\newtheorem{lemma}[theorem]{Lemma}
\newtheorem{proposition}[theorem]{Proposition}
\newtheorem{corollary}[theorem]{Corollary}
\theoremstyle{remark}\newtheorem{remark}[theorem]{Remark}
\newcommand{\R}{\mathbb R}\newcommand{\C}{\mathbb C}\newcommand{\N}{\mathbb N}
\newcommand{\eps}{\varepsilon}\newcommand{\HH}{\mathcal H}\newcommand{\one}{\mathbf1}
\newcommand{\DS}{\mathrm{DSASM}}\newcommand{\ASM}{\mathrm{ASM}}
\newcommand{\Sone}{\mathcal S_1}\newcommand{\Stwo}{\mathcal S_2}
\newcommand{\ceil}[1]{\left\lceil#1\right\rceil}
\DeclareMathOperator{\tr}{tr}\DeclareMathOperator{\Ran}{Ran}
\DeclareMathOperator{\Log}{Log}\DeclareMathOperator{\Var}{Var}
\DeclareMathOperator{\Cov}{Cov}\DeclareMathOperator{\diag}{diag}
\hypersetup{pdftitle={Report 128 Full leading equivalent for diagonally symmetric alternating sign matrices},pdfauthor={}}
\begin{document}
\begin{center}
{\small REPORT 128}\par\vspace{8pt}
{\LARGE\bfseries The full leading equivalent for\par}
\vspace{4pt}
{\LARGE\bfseries diagonally symmetric\par}
\vspace{4pt}
{\LARGE\bfseries alternating sign matrices\par}
\vspace{9pt}
{\large Individual fugacity amplitudes and finite-section boundary limits}\par
\vspace{9pt}{\small 2 October 2026}\par
{\small A proof with an unchanged earlier-input appendix and offline replay}
\end{center}
\begin{abstract}
For the number $a_n$ of diagonally symmetric alternating sign matrices,
A005163, we prove a full-sequence equivalent
$a_n\sim C e^{\alpha n^2+\beta n}n^{-5/72}$, where
$\alpha=\tfrac12\log(3\sqrt3/4)$ and $\beta=\tfrac14\log3$.
The positive constant is characterized by convergent Fredholm determinants
and a convergent scalar Fourier series. More generally, the partition
function at every positive diagonal fugacity has a common individual
amplitude, with a locally holomorphic limit on the right half-plane.
There are two distinct new steps: convergence of the companion determinant
remainder, and a quantitative finite-section boundary theorem that removes
the alternating ambiguity in adjacent products. The latter retains the
leading Galerkin correction and uses an exact Christoffel norm identity
for weighted continuous-Hahn tightness. We give the amplitude-aware inverse
threshold with a vanishing remainder before integer ceilings, and the
constant terms of all fixed diagonal cumulants. No numerical amplitude,
explicit rate of convergence, inverse-power correction to the general
count, or all-orders expansion is claimed.
\end{abstract}
\newpage
\setcounter{tocdepth}{1}
{\small\setlength{\parskip}{0pt}\tableofcontents}
\newpage
\section{Definitions, theorem and proof dependencies}\label{sec:main}
An alternating sign matrix has entries in $\{-1,0,1\}$, alternating
nonzero entries along each row and column, and all row and column sums
one. It is diagonally symmetric when it equals its transpose. Write
\[
 a_n=|\DS(n)|,\qquad S(A)=|\{i:A_{ii}\ne0\}|,\qquad
 Z_n(s)=\sum_{A\in\DS(n)}s^{S(A)},\qquad Z_0(s)=1.
\]
Thus $a_n=Z_n(1)$ and $a_n$ is A005163. All logarithms of positive
numbers are natural. Put
\begin{equation}\label{constants}
 c=\sqrt3,\quad \alpha=\tfrac12\log\frac{3\sqrt3}{4},\quad
 \beta=\tfrac14\log3,\quad \kappa=\tfrac5{72},\quad
 b=\tfrac34\log3-\tfrac12\log2.
\end{equation}
On $\HH_+=\{\Re s>0\}$ define
\begin{equation}\label{pressure}
 P(s)=\frac{(s+1)^2}{2(s+2)},\qquad
 B(s)=b+\tfrac12\Log P(s).
\end{equation}
Here $\Log P(s)=2\Log(s+1)-\Log(s+2)-\log2$, with the indicated
right-half-plane logarithms. In particular $P(c)=1$ and $B(1)=\beta$.
The companion polynomial and its centered logarithm are
\begin{equation}\label{companiondefs}
 D_n(t)=\frac{Z_n(\sqrt t)Z_{n+1}(\sqrt t)}{\sqrt t},\quad
 f(t)=\Log P(\sqrt t),\quad
 E_n(t)=\Log\frac{D_n(t)}{D_n(3)}-nf(t),
\end{equation}
where $t\in\Omega_t=\C\setminus(-\infty,0]$ and the square root
has positive real part. The quotient defining $D_n$ is interpreted
polynomially. Its logarithm is normalized to vanish at $t=3$.

\begin{theorem}[Full leading equivalent]\label{thm:main}
There is a holomorphic function $E$ on $\Omega_t$, real on $(0,\infty)$,
with $E(3)=0$, such that $E_n\to E$ locally uniformly. Define
\begin{equation}\label{Ccal}
 C_{\rm cal}=
 \frac{3^{11/72}\pi^{1/6}e^{\zeta'(-1)/6}}
      {2^{1/24}\Gamma(1/3)^{1/3}}>0
\end{equation}
and, on $\HH_+$,
\begin{equation}\label{amplitude}
 C(s)=C_{\rm cal}\exp\{E(s^2)/2\}
       (s/c)^{1/2}P(s)^{-1/4}.
\end{equation}
The powers use the analytic logarithms normalized at $s=c$.
Then $C$ is nowhere zero and positive for $s>0$, and
\begin{equation}\label{fullequivalent}
 Z_n(s)=C(s)e^{\alpha n^2+nB(s)}n^{-\kappa}(1+o(1))
\end{equation}
locally uniformly on $\HH_+$, along the full integer sequence.
In particular,
\begin{equation}\label{aequiv}
 a_n\sim C_{\rm cal}2^{-1/4}e^{E(1)/2}
               e^{\alpha n^2+\beta n}n^{-5/72}.
\end{equation}
A convergent operator/Fourier expression for $E(t)$ is given in
\eqref{Eexplicit}; it defines the constant without a numerical fit.
\end{theorem}

\subsection{What is supplied in the reader bundle}
The unchanged earlier article \cite{r124}, supplied both as a PDF and
as its complete TeX source in the bundle's \texttt{earlier\_input}
folder, proves the bounded-remainder stage. Its scope statements are
historical: its failure to assert a limit is strengthened here. All
new convergence, calibration and boundary steps are proved below;
no unpublished proof note is a required reader input.

For precision, these are the earlier results used, with their locations
in that supplied article:
\begin{enumerate}[label=\textup{(I\arabic*)}]
\item Sections 2--4: imaginary-axis zeros of $Z_n$, negative-real zeros
of $D_n$, cross-size interlacing, the exact adjacent determinant and
its finite Pascal reduction, and
$Z_n(c)Z_{n+1}(c)=c2^n A_{n+1}$, where $A_N=|\ASM(N)|$.
\item Sections 6--8: the exact degree-space compression
$D_n(t)/D_n(3)=\det(I+(t-3)P_n\mathcal B P_n)$, the bounded operator
$\mathcal B$, its unweighted decomposition $C_0+\mathcal E$ with
$0\le C_0\le I/3$, $\mathcal E\in\Stwo$,
$\mathcal S=(\mathcal E+\mathcal E^*)/2\in\Sone$, and the locally
uniform nonzero relative determinant $\mathcal A(t)$.
The exact continuous-Hahn/Fourier/Jacobi transform and the gamma-phase
comparison are included there.
\item Sections 9--11: the normalized Jacobi expansion, its endpoint
and derivative bounds, the mean formula
$\tr(J_nFJ_n)=(n+1)\pi^{-1}\int_0^\pi F+o(1)$ for endpoint-vanishing
$F$ with $\int |F|/\min(\theta,\pi-\theta)<\infty$, the global
angular kernel bound $|K_n(\theta,\phi)|\le C/|\theta-\phi|$, the
finite tilted-projection variance bound, and the Hilbert--Schmidt
weighted Hardy wave comparisons.
\item Sections 12 and 14: the convergent trace-class cross factor
$\Xi(t)$, the compact-uniform scalar bound
$\log T_n(t-3)-2nf(t)=O_K(1)$, the analogous companion bound, and
local uniform bounds on adjacent individual logarithmic differences.
\end{enumerate}
These are theorems and exact identities proved in \cite{r124}, rather
than assertions supplied by finite checks. We restate the definitions
needed to apply them. The additional published ingredients are the
ASM/OSASM product asymptotics in \cite{bfk}, the continuous-Hahn
normalization and generating function \cite{cll,dlmf}, and the
\emph{polynomial} Jacobi-statistic central limit theorem \cite{bd}.
We state precisely where each is used and prove the necessary extensions.

\subsection{The two logically separate limits}
An adjacent-product limit does not alone determine the individual
amplitudes: a residual $L+(-1)^n d$ has a constant adjacent sum.
Section~\ref{sec:companion} first proves $E_n\to E$.
Sections~\ref{sec:boundary}--\ref{sec:edge} then prove the distinct
boundary-ratio limit
\begin{equation}\label{Ftarget}
 F_n(s):=\frac{sZ_n(s)/Z_{n+1}(s)}{cZ_n(c)/Z_{n+1}(c)}
 \longrightarrow \frac{s/c}{\sqrt{P(s)}}.
\end{equation}
The finite-section proof is at the $1/n$ coefficient scale. In
particular, it does not substitute an infinite resolvent for a finite
one merely from strong convergence. Section~\ref{sec:consequences}
combines the two limits, proves the inverse threshold and derives
constant-order cumulants.

\section{Convergence of the companion remainder}\label{sec:companion}
We use $\delta=t-3\in\Omega=\C\setminus(-\infty,-3]$.
All analytic logarithms in this section vanish at $\delta=0$.
Distinguish the original real-line weight $w_y$ in
Section~\ref{sec:boundary} from the angular multiplier
\begin{equation}\label{wang}
 w(\theta)=\frac1{1+2\cosh(\xi(\theta)/3)},\quad
 \xi(\theta)=2\log\cot(\theta/2),\quad s_0(\theta)=\sqrt{w(\theta)}.
\end{equation}
Thus $0\le w\le1/3$ and $w=O(d^{2/3})$ for
$d(\theta)=\min(\theta,\pi-\theta)$.
Let $U:\ell^2(\N_0)\to L^2(0,\pi)$ synthesize the real normalized
Jacobi angular basis with parameters $(4/3,2/3)$, and let $H$ be the
orthogonal Hardy projection obtained from the negative real half-line
under Fourier and angular transformation. Explicitly it is unitarily
conjugate to the Fourier image of multiplication by $\one_{(-\infty,0)}$.
Write $P_n$ for the first $n$ coordinate projection and $J_n=UP_nU^*$.
Set
\[
 X=M_{s_0}HU,\quad Y=M_{s_0}(I-H)U,\quad
 A_n=P_nX^*XP_n,\quad B_n=P_nY^*YP_n.
\]
Here $A_n,B_n$ act on $\Ran P_n$. Define
\[
 Q_n(\delta)=\det(I+\delta A_n)=\det(I+\delta B_n),\quad
 S_n(\delta)=\det(I+\delta(A_n+B_n)),
\]
\[
 T_n(\delta)=\det_{\Ran J_n}(I+\delta J_nM_wJ_n).
\]
Equality of the two single determinants follows from complex
conjugation in the real Jacobi basis, which interchanges $H$ and $I-H$.
The supplied earlier proof gives, locally uniformly on $\Omega$,
\begin{equation}\label{inheritedfactors}
 \frac{D_n(t)/D_n(3)}{Q_n(t-3)}\longrightarrow\mathcal A(t),
 \qquad \frac{S_n(t-3)}{T_n(t-3)}\longrightarrow\Xi(t).
\end{equation}
For clarity, their complete operator definitions are
\begin{align}
 \mathcal A(t)&=\det{}_2\!\left(I+\delta(I+\delta C_0)^{-1}\mathcal E\right)
 \exp\!\left\{\delta\tr[(I+\delta C_0)^{-1}\mathcal S]\right\},
 \label{Arel}\\
 \Xi(t)&=\det\!\left[(I+\delta\mathcal D)(I+\delta M_w)^{-1}\right],
 \quad \mathcal D=HM_wH+(I-H)M_w(I-H).
 \label{Xidef}
\end{align}
In \eqref{Arel} the real-line operators are exactly those defined in
\eqref{kernelreal} below. The difference $\mathcal D-M_w$ is trace
class; the two factors are invertible on $\Omega$. Both $\mathcal A$
and $\Xi$ are nowhere zero, normalized to one at calibration, and
positive on $t>0$. The gamma-phase comparison has limiting relative
determinant one, which is already included in the first limit in
\eqref{inheritedfactors}.

\subsection{A moving-edge mixed determinant limit}
The weighted-wave result of \cite[Section 11]{r124} gives bounded maps
\[
 X_0e_j=\frac{s_0(\theta)e^{-i(j+3/2)\theta+ic_0}}{\sqrt{2\pi}},
 \qquad Y_0e_j=\frac{s_0(\theta)e^{i(j+3/2)\theta-ic_0}}{\sqrt{2\pi}},
 \quad c_0=11\pi/12,
\]
with $X-X_0,Y-Y_0\in\Stwo$. Importantly, $XY^*=0$ exactly,
since $UU^*=I$ and $H(I-H)=0$.
Define $M_n$ to be multiplication by $e^{-in\theta}$ and define the
Hilbert--Schmidt operator $K$ by its kernel
\begin{equation}\label{Kedge}
 K(\theta,\phi)=\frac{s_0(\theta)s_0(\phi)}{2\pi}
 \frac{e^{2ic_0-(3i/2)(\theta+\phi)}}{1-e^{-i(\theta+\phi)}}.
\end{equation}
At either singular corner its square is bounded by
$C\theta^{2/3}\phi^{2/3}/(\theta+\phi)^2$ or its reflection,
which is integrable after setting $\theta=ru$, $\phi=r(1-u)$.
Thus $K\in\Stwo$.

The geometric sum gives the exact identity
\[
 X_0P_nY_0^*=K-M_nKM_n.
\]
Since $M_n\rightharpoonup0$ by the Riemann--Lebesgue lemma and $K$
is compact, $M_nKM_n\to0$ strongly. Taking the strong limit gives
$X_0Y_0^*=K$. Set $E_X=X-X_0$, $E_Y=Y-Y_0$. Finite-rank approximation
in $\Stwo$ gives
\[
 XP_nY^*-X_0P_nY_0^*=E_XP_nY^*+X_0P_nE_Y^*
       \longrightarrow E_XY^*+X_0E_Y^*=-K
\]
in Hilbert--Schmidt norm. Consequently
\begin{equation}\label{movingK}
 K_n:=XP_nY^*=-M_nKM_n+o_{\Stwo}(1).
\end{equation}
In particular $\tr(A_nB_n)=\|K_n\|_2^2\to\|K\|_2^2$.
The moving term in \eqref{movingK} is retained; it does not tend to
zero in Hilbert--Schmidt norm.

Let $E_I$ extend functions on $(0,\pi)$ by zero to the circle
$(0,2\pi)$. Let $P_+^\circ$ project onto circle modes $j\ge1$ and
$P_-^\circ$ onto $j\le-1$, and set $G_\pm=M_{e^{\pm3i\theta/2}}$.
Define
\begin{equation}\label{edgeAB}
 A_e=M_{s_0}G_-E_I^*P_+^\circ E_IG_+M_{s_0},\qquad
 B_e=M_{s_0}G_+E_I^*P_-^\circ E_IG_-M_{s_0}.
\end{equation}
The strict mode cutoffs follow from $n-j=1,\ldots,n$ in the finite
wave sums, and yield
\begin{equation}\label{diagonallimits}
 M_n^*XP_nX^*M_n\longrightarrow A_e,\qquad
 M_nYP_nY^*M_n^*\longrightarrow B_e
\end{equation}
strongly. Here is why the fixed Hilbert--Schmidt errors do not affect
these limits. For fixed $h$, $P_nX^*M_nh$ is bounded and weakly null:
pairing with $g$ reduces to pairing $M_nh$ with $XP_ng\to Xg$ in
norm. Compactness of $E_X$ then gives
$E_XP_nX^*M_nh\to0$, and compactness also gives $E_X^*M_nh\to0$.
Expand $XP_nX^*-X_0P_nX_0^*=E_XP_nX^*+X_0P_nE_X^*$ to obtain
the first claim. The other is identical with reversed modulation.
Both limits are positive and bounded by $I/3$.

For $0\le T\le I/3$, $(I+\delta T)^{-1}$ is uniformly bounded
for $\delta$ in compact subsets of $\Omega$. Strong convergence
in \eqref{diagonallimits} gives strong convergence of these resolvents,
locally uniformly on each fixed vector. A finite parameter net and
the resolvent identity justify this uniformity.

Apply Sylvester's determinant identity to the block map with entries
$XP_n$ and $YP_n$, and take its upper-left Schur complement. With
\[
 R_{X,n}=(I+\delta XP_nX^*)^{-1},\qquad
 R_{Y,n}=(I+\delta YP_nY^*)^{-1},
\]
one obtains, with the order of factors as written,
\begin{equation}\label{schur}
 \frac{S_n(\delta)}{Q_n(\delta)^2}
 =\det[I-\delta^2R_{Y,n}K_n^*R_{X,n}K_n].
\end{equation}
Conjugate the product on the right by $M_n$, and use
\eqref{movingK}. Products of two Hilbert--Schmidt factors are trace
class, so all terms containing its $o_{\Stwo}(1)$ error tend to zero
in trace norm, uniformly on compact parameter sets. The main term is
\[
 \delta^2(M_nR_{Y,n}M_n^*)K^*(M_n^*R_{X,n}M_n)K.
\]
Strong resolvent convergence and finite-rank approximation of $K$
now give locally uniform trace-norm convergence. Hence
\begin{equation}\label{Psidef}
 \frac{S_n(\delta)}{Q_n(\delta)^2}\longrightarrow\Psi(\delta):=
 \det\!\left[I-\delta^2(I+\delta B_e)^{-1}K^*
                              (I+\delta A_e)^{-1}K\right].
\end{equation}
Every finite ratio is nonzero on $\Omega$, and $\Psi(0)=1$.
Hurwitz's theorem makes $\Psi$ nowhere zero. It is positive for
real $\delta>-3$, as the nonzero limit of positive ratios. Thus
\begin{equation}\label{mixedlimit}
 2\log Q_n(\delta)-\log S_n(\delta)\longrightarrow-\Log\Psi(\delta)
\end{equation}
locally uniformly, with the common normalization at zero. Alternatively,
the modulated positive block matrices have strong limit
$\bigl(\begin{smallmatrix}A_e&-K\\-K^*&B_e\end{smallmatrix}\bigr)$
bounded by $I/3$; its Schur complement directly proves invertibility
of the limiting factor in \eqref{Psidef}.

\subsection{The scalar constant: a polynomial theorem and its extension}
For an angular function $F(\theta)=g(\cos\theta)$ put
\[
 \mu(F)=\frac1\pi\int_0^\pi F(\theta)\,d\theta,\quad
 F_k=\frac1\pi\int_0^\pi F(\theta)\cos(k\theta)\,d\theta,\quad
 v(F)=\sum_{k\ge1}kF_k^2.
\]
For real $F$, the squares in $v$ are ordinary real squares.

\begin{proposition}\label{prop:scalarconstant}
If $F$ is real and angular H\"older of exponent
$1/2<\gamma\le1$, with $F(0)=F(\pi)=0$, then $v(F)<\infty$ and
\begin{equation}\label{scalarconstant}
 \log\det_{\Ran J_n}(J_ne^gJ_n)
     =(n+1)\mu(F)+\tfrac12v(F)+o(1).
\end{equation}
Here $g$ denotes multiplication by the function whose angular form is $F$.
\end{proposition}
\begin{proof}
The rank-$n$ Jacobi polynomial ensemble has density proportional to
its squared Vandermonde times the Jacobi weight. Andreief's finite
identity gives
\[
 \mathbb E e^{X_n(g)}=\det(J_ne^gJ_n),\quad
 \mathbb E X_n(g)=\tr(J_ngJ_n),\quad
 \Var_n(g)=\|(I-J_n)gJ_n\|_2^2.
\]
The earlier global kernel bound gives for every real angular
$C^\eta$ function $H$, $\eta>1/2$,
\begin{equation}\label{varianceholder}
 \Var_n(H)=\tfrac12\iint(H(\theta)-H(\phi))^2
                      |K_n(\theta,\phi)|^2\,d\theta\,d\phi
 \le C_\eta[H]_{C^\eta}^2,
\end{equation}
independently of $n$. Endpoint vanishing is unnecessary for this bound.

The polynomial input is Corollary 2.2 of Breuer--Duits
\cite{bd}: if the recurrence matrix has a Laurent right limit,
every fixed real polynomial statistic, centered by its exact mean,
converges to a centered Gaussian with variance
$\sum_{k\ge1}k\widehat p_k\widehat p_{-k}$, where the Fourier
coefficients are those of the polynomial applied to the Laurent symbol.
The fixed Jacobi recurrence here has off-diagonal limit $1/2$ and
diagonal limit zero, hence the unique full-sequence symbol
$(z+z^{-1})/2$. Its coefficients are bounded and the support is compact.
Thus this theorem applies to every real polynomial $p$, with variance
$v(p(\cos\theta))$. Only the polynomial theorem is used; its separate
$C^1$ extension is not applied to the fractional endpoint symbol.

A weak CLT alone is insufficient for the determinant. To obtain
exponential integrability, let $E_r=e^{rg/2}$ and
$T_r=J_ne^{rg}J_n$ on $\Ran J_n$. The isometry
$E_rJ_nT_r^{-1/2}$ has range projection
$Q_r=E_rJ_nT_r^{-1}J_nE_r$. For any bounded real multiplier $a$,
commutation with $E_r$ gives the exact identity
\[
 (I-Q_r)aE_rJ_nT_r^{-1/2}
   =(I-Q_r)E_r(I-J_n)aJ_nT_r^{-1/2}.
\]
For $r\ge0$, taking Hilbert--Schmidt norms gives
$\Var_{rg}(a)\le e^{r\operatorname{osc}(g)}\Var_n(a)$.
Taylor's integral formula for the logarithmic moment generating
function, also applied to $-g$, therefore yields for real $u$
\begin{equation}\label{exponentialbound}
 \log\mathbb E e^{u(X_n(g)-\mathbb EX_n(g))}
 \le \tfrac12u^2e^{|u|\operatorname{osc}(g)}\Var_n(g).
\end{equation}
For fixed polynomial $p$, the bound at $u=2$ and
\eqref{varianceholder} prove uniform integrability of
$e^{X_n(p)-\mathbb EX_n(p)}$. The polynomial CLT consequently implies
\begin{equation}\label{polynomialG}
 G_n(p):=\log\mathbb E e^{X_n(p)-\mathbb EX_n(p)}\longrightarrow v(p)/2.
\end{equation}

We next justify extension uniformly in $n$. For real $g,h$ put
$a=g-h$ and $g_u=h+ua$, $0\le u\le1$. Differentiation of the
finite determinant and then integration along its tilt path give
\[
 \frac{d}{du}G_n(g_u)
  =\mathbb E_{g_u}X_n(a)-\mathbb E_0X_n(a)
  =\int_0^1\Cov_{rg_u}(X_n(a),X_n(g_u))\,dr.
\]
Tilted covariance Cauchy--Schwarz and the preceding ordered identity
show
\begin{equation}\label{Gcontinuity}
 \left|\frac{d}{du}G_n(g_u)\right|
 \le e^{\operatorname{osc}(g_u)}
       \sqrt{\Var_n(a)\Var_n(g_u)}.
\end{equation}
By \eqref{varianceholder}, $G_n$ are uniformly Lipschitz on bounded
real $C^\eta$ sets for $\eta>1/2$.

Choose $1/2<\eta<\gamma$. Extend $F$ evenly to the circle.
Its Fej\'er convolutions $F_m$ are even trigonometric polynomials,
converge uniformly, and have bounded $C^\gamma$ seminorms.
The elementary increment estimate
\[
 [H]_{C^\eta}\le C\|H\|_\infty^{1-\eta/\gamma}
                        [H]_{C^\gamma}^{\eta/\gamma}
\]
proves convergence in $C^\eta$. Subtract
$a_m+b_m\cos\theta$ with
$a_m=(F_m(0)+F_m(\pi))/2$ and
$b_m=(F_m(0)-F_m(\pi))/2$. The corrected functions $Q_m$ vanish
at both endpoints, are polynomials in $\cos\theta$, and converge
to $F$ in $C^\eta$, since $a_m,b_m\to0$.

The circle $H^{1/2}$ difference-quotient identity, first verified by
Fourier orthogonality for trigonometric polynomials, identifies
$\sum |k||\widehat H_k|^2$ with a fixed positive multiple of
\[
 \iint\frac{|H(\theta)-H(\phi)|^2}
                   {\sin^2((\theta-\phi)/2)}\,d\theta\,d\phi.
\]
It is bounded by $C_\eta[H]_{C^\eta}^2$. Thus $v(F-Q_m)\to0$ and
$v(Q_m)\to v(F)$ by weighted Cauchy--Schwarz. Equations
\eqref{polynomialG} and \eqref{Gcontinuity}, first with $n\to\infty$
and then $m\to\infty$, give $G_n(g)\to v(F)/2$.
Finally the earlier Jacobi mean theorem applies because
$F=O(d^\gamma)$ at the endpoints. It gives
$\mathbb EX_n(g)=(n+1)\mu(F)+o(1)$, proving
\eqref{scalarconstant}.
\end{proof}

Apply this proposition to
\[
 F_\delta(\theta)=\Log(1+\delta w(\theta)),\qquad
 F_{\delta,k}=\frac1\pi\int_0^\pi F_\delta(\theta)\cos(k\theta)\,d\theta.
\]
For real $\delta>-3$ it is real, endpoint-vanishing and $C^{2/3}$,
and the exact pressure integral from \cite[Section 9]{r124} is
$\mu(F_\delta)=2f(3+\delta)$. Therefore
\begin{equation}\label{scalarreallimit}
 \log T_n(\delta)-2nf(3+\delta)
 \longrightarrow 2f(3+\delta)+\tfrac12\sum_{k\ge1}kF_{\delta,k}^2.
\end{equation}
The earlier complex scalar bound makes the left side locally bounded
and holomorphic on $\Omega$. Vitali's theorem extends this convergence
locally uniformly over $\Omega$. The displayed series is its analytic
expression there, with squares, not absolute squares. Indeed for
$1/2<\eta<2/3$ the map $\delta\mapsto F_\delta$ is holomorphic
as a $C^\eta$-valued map, by the uniform lower bound for
$|1+\delta u|$, $0\le u\le1/3$, on compact parameter sets.
Its image on a compact set is compact in the weighted Fourier
$\ell^2$ space. Finite-coordinate projections then have uniformly
vanishing tails on this image. The bilinear Cauchy--Schwarz inequality
proves local uniform convergence and holomorphy of the square series.
The real identity and analytic uniqueness prove the claimed expression.

\subsection{The explicit limit and branch normalization}
Combine \eqref{inheritedfactors}, \eqref{mixedlimit}, and
\eqref{scalarreallimit}. Every constituent logarithm is normalized at
$t=3$, and their locally uniform limits give
\begin{equation}\label{Eexplicit}
 \boxed{\ E(t)=\Log\mathcal A(t)+\tfrac12\Log\Xi(t)
       -\tfrac12\Log\Psi(t-3)+f(t)
       +\tfrac14\sum_{k\ge1}kF_{t-3,k}^2.\ }
\end{equation}
This proves $E_n\to E$ locally uniformly on $\Omega_t$.
At $t=3$ every term is zero. For positive $t$ the three determinant
factors are positive, so $E(t)$ is finite and real. Notice both the
minus sign of the mixed factor and the scalar coefficient $1/4$:
one first solves $S_n/Q_n^2\to\Psi$ for $\log Q_n$, and then
halves the scalar variance term. These fix the amplitude normalization.
As a separate normalization check, for the polynomial $p(x)=x$ the
only positive cosine coefficient is $1/2$, so $v(p)=1/4$ and the
centered logarithmic determinant tends to $1/8$. Directly the finite
variance is the square of the single Jacobi recurrence coefficient
crossing the degree-$n$ boundary, which tends to $1/4$.

\section{Calibration and the exact boundary resolvent}\label{sec:boundary}
\subsection{A common calibration amplitude on both parities}
The earlier exact identities give
\[
 Z_{2m}(c)=2^m\prod_{k=1}^m\frac{A_{2k}}{A_{2k-1}},\qquad
 Z_n(c)=c2^{n-1}A_n/Z_{n-1}(c)\quad(n\text{ odd}).
\]
Let $O_m$ denote the number of off-diagonally symmetric ASMs of
order $2m$. The classical products in \cite{bfk} are
\[
 A_N=\prod_{i=0}^{N-1}\frac{(3i+1)!}{(N+i)!},\qquad
 O_m=\prod_{i=1}^{m}\frac{(6i-2)!}{(2m+2i)!}.
\]
Successive ratios give
\[
 \frac{O_m}{O_{m-1}}
 =\frac{(6m-2)!(2m)!}{(4m)!(4m-2)!}
 =\frac12\frac{A_{2m}}{A_{2m-1}}.
\]
For example, the old denominator factors in the $O_m$ quotient
telescope to $(4m-2)!/(2m)!$. The initial products equal one,
so $Z_{2m}(c)=2^{2m}O_m$ exactly.

The published product asymptotics \cite[Equations (11.3)--(11.6)]{bfk},
using the even-order OSASM branch, state
\begin{align*}
 A_n&\sim C_{\rm ASM}e^{2\alpha n^2}n^{-2\kappa},\\
 2^nO_{n/2}&\sim C_{\rm cal}e^{\alpha n^2+bn}n^{-\kappa}
                    \quad(n\text{ even}),\\
 C_{\rm ASM}&=
 \frac{2^{5/12}\pi^{1/3}e^{\zeta'(-1)/3}}
      {3^{7/36}\Gamma(1/3)^{2/3}}.
\end{align*}
These are known product asymptotics, not an assumption about DSASM
amplitudes. Substitution in the odd exact identity gives the odd
constant
\[
 C_{\rm odd}=(c/2)e^{b-\alpha}C_{\rm ASM}/C_{\rm cal}.
\]
Now $b-\alpha=\tfrac12\log2$ and
$C_{\rm ASM}/C_{\rm cal}^2=\sqrt{2/3}$ by the displayed constants.
Therefore $C_{\rm odd}=C_{\rm cal}$, proving the full-sequence
calibration
\begin{equation}\label{calibrationequiv}
 Z_n(c)\sim C_{\rm cal}e^{\alpha n^2+bn}n^{-\kappa}.
\end{equation}
Only this equivalent is needed below; no finite-size correction is
transferred to a noncalibrated fugacity.

\subsection{The fixed operator, basis and source}
Put $h=\sqrt3$, $a=\pi/(3h)$, $p(y)=y^2+3/4$, and
\begin{equation}\label{realdefinitions}
 w_y(y)=\frac{p(y)}{2(2\cosh(2ay)-1)},\quad
 m(y)=\frac1{1+e^{2ay}},\quad
 \varphi(u)=\frac1{3(1+2\cosh(2au))}.
\end{equation}
Let $\psi_j$ be the real orthonormal polynomials in
$\mathcal H=L^2(w_y\,dy)$ whose leading coefficients are
$(-1)^j/\sqrt{H_j}$. The degree-$<n$ projection is $P_n$.
On polynomials put
\[
 (Tg)(y)=\int_\R e^{au}\varphi(u)g(y+u)\,du,\qquad \mathcal B=mT.
\]
The boundedness of $\mathcal B$ is part of the earlier proof.
The unweighted transform of $\mathcal B$ has kernel
\begin{equation}\label{kernelreal}
 \widetilde{\mathcal B}(y,x)=g_0(y)\varphi(x-y)h_0(x),\quad
 g_0=\sqrt{w_y}\,m e^{-ay},\quad h_0=e^{ay}/\sqrt{w_y},
 \quad C_0(y,x)=\one_{y<0}\varphi(x-y)\one_{x<0}.
\end{equation}
Here $\mathcal E=\widetilde{\mathcal B}-C_0$ is the operator in
\eqref{Arel}. The polynomial inverse of $T$ is
\[
 (\mathcal Kg)(y)=g(y)+e^{i\pi/3}g(y+ih)+e^{-i\pi/3}g(y-ih).
\]
The exact degree-space determinant is
$D_n(t)/D_n(3)=\det(I+(t-3)P_n\mathcal B P_n)$.

Define $d_j=\sqrt{(j+1)(j+2)}$, $h_j=j!(3)_j$, and
\begin{equation}\label{Hnorm}
 H_j=2(27/16)^j
       \frac{j!(3)_j(5/3)_j(7/3)_j}{(3/2)_j(5/2)_j}.
\end{equation}
The affine source $q_t\in\mathcal H$ is
\begin{equation}\label{qsource}
 q_t(y)=\frac{A_t e^{2ay}+(-y-t+3/2)e^{ay}/\sqrt2+C_t
           +[(1-t)y/2-t+3/2]e^{-ay}/\sqrt2}{p(y)\cosh(ay)},
\end{equation}
where $A_t=\sqrt{3/2}(t+1)/4$ and $C_t=\sqrt{3/2}(t-2)$.
Write $q_t=q_3+(t-3)q^{(1)}$ and
\begin{equation}\label{gamma}
 \gamma_j(t)=\left\langle\psi_j,
   (I+(t-3)P_{j+1}\mathcal B^*P_{j+1})^{-1}P_{j+1}q_t\right\rangle,
 \qquad a_j^{\rm src}=\langle\psi_j,q_3\rangle.
\end{equation}
The inner product is linear in its second argument. The inverse in
\eqref{gamma} acts on the finite-dimensional range, never on an
unjustifiably substituted infinite section.

\begin{proposition}[Exact all-parity boundary formula]\label{prop:boundary}
For every $n\ge1$ and $s>0$, with $t=s^2$,
\begin{equation}\label{boundaryexact}
 \frac{sZ_n(s)}{Z_{n+1}(s)}
 =(-1)^{n-1}d_{n-1}\sqrt{h_{n-1}/H_{n-1}}\,\gamma_{n-1}(t).
\end{equation}
It follows that $F_n(s)=\gamma_{n-1}(s^2)/a_{n-1}^{\rm src}$, and
\begin{equation}\label{calcoefficient}
 a_j^{\rm src}\sim\frac{(-1)^j\sqrt2\,3^{-1/4}}{j+1}.
\end{equation}
\end{proposition}
\begin{proof}
We supply the borders and source normalization explicitly.
Let
\[
 K_t(u,v)=\frac{v-u}{1-uv}
       \left(t+\frac{(1+u)(1+v)}{1-u-v}\right),\qquad
 M_{ij}=[u^iv^j]K_t(u,v).
\]
The source Pfaffian formula \cite[Equation (4.1)]{bfk}
and the earlier adjacent determinant give
$H_n=(M_{i,j+1})_{0\le i,j<n}$ with $\det H_n=D_n(t)$.
Set $b_i=M_{i0}$ and $v_t=te_0-b$; directly
$b_0=0$, $b_1=-(t+1)$ and $b_i=-2$ for $i\ge2$.
For even $n$, replacing the last column of $H_n$ by $b$ and cycling
it to the front gives
$-\det(H_n)e_{n-1}^{\mathsf T}H_n^{-1}b=Z_n(s)^2$.
The $e_0$ cofactor vanishes because its complementary skew matrix
has odd size. For odd $n$ the full skew determinant vanishes, while
the $e_0$ complementary skew minor indexed $1,\ldots,n-1$ has
determinant $(Z_n(s)/s)^2$. Including the cofactor signs in either
case proves
\begin{equation}\label{cofactor}
 sZ_n(s)/Z_{n+1}(s)=e_{n-1}^{\mathsf T}H_n^{-1}v_t.
\end{equation}

Use the finite matrices $S$ (lower shift), $L_{ij}=\binom ij$,
$E_{ii}=(-1)^i$, $\Pi_{ij}=\binom{i+j}{i}$, and $T_n(g)=g(S)$.
The earlier finite Pascal factorization, or direct multiplication
of the displayed generating functions, is
\[
 H_n=T_n(B_0)LEJ_n(t)EL^{\mathsf T},\quad
 J_n(t)=I+T_n(G_t)\Pi,
\]
where
\[
 B_0(z)=\frac{(z-2)(z+1)(2z-1)}{(1-z)(z^2-z+1)},\quad
 G_t(z)=\frac{(t-4)z^2+(4-t)z+t-1}{(2-z)(1-z)^2(1+z)}.
\]
The last row of $L^{-\mathsf T}$ is $e_{n-1}^{\mathsf T}$.
The generating function of $v_t$ is $V_t(z)=t+(t+1)z+2z^2/(1-z)$,
so the binomial transform
$(1-z)^{-1}[V_t/B_0](-z/(1-z))$ simplifies to
\[
 w_t(z)=\frac{(t-z)(z^2-z+1)}{(2-z)(1-z)^2(1+z)}.
\]
Consequently \eqref{cofactor} equals
$(-1)^{n-1}e_{n-1}^{\mathsf T}J_n(t)^{-1}w_t^{(n)}$,
where the vector contains its first $n$ coefficients.

Put $D=\diag(d_j)$, $A=D^{-1}SD$ and
$C=D^{-1}VD$ with $V_{ij}=\binom{i+j+2}{i}$.
For $r(z)=1-z+z^2$, $q(z)=2+z-z^2$, the exact finite identity is
$D^{-1}J_nD=I+C+(t-3)[r(A)/q(A)]C$.
Multiplying by $q(A)/r(A)$ gives the old row-action Gram matrix
$F_n^\phi=\mathcal K_n^{\rm row}(I+C)+(t-3)C$.
The forcing is especially simple:
\[
 (q/r)w_t=(t-z)/(1-z)^2,\qquad
 p_i(t)=[t(i+1)-i]/d_i.
\]
Thus \eqref{cofactor} becomes
$(-1)^{n-1}d_{n-1}e_{n-1}^{\mathsf T}(F_n^\phi)^{-1}p(t)$.

Let $\phi_j=u_j/\sqrt{h_j}$ be the Meixner--Pollaczek basis from
\cite[Section 6]{r124}, with
\[
 \sum_{j\ge0}u_j(y)z^j/j!=r(z)^{-3/2}e^{y\eta(z)},\qquad
 \eta'(z)=-1/r(z),\quad\eta(0)=0.
\]
The lower-triangular basis change $\psi=U_n\phi$ has last diagonal
entry $\sqrt{h_{n-1}/H_{n-1}}$. Congruence gives
$F_n^\psi=U_nF_n^\phi U_n^{\mathsf T}=\mathcal K_n^{\rm row}
+(t-3)M_n$, where $(M_n)_{ij}=\langle\psi_i,m\psi_j\rangle$.
The matrix $(\mathcal K_n^{\rm row})^{-1}M_n$ is exactly the matrix
of $P_n\mathcal B^*P_n$ in column coordinates. If
$\ell_t(\phi_i)=p_i(t)$, its transformed source has entries
$\ell_t(T\psi_i)$. We now identify their genuine representing vector.

Define the exponentially decaying density
\[
 \nu_t(y)=\frac{(y+t-3/2)e^{ay}+[(1-t)y+(3-t)/2]e^{-ay}}
                    {2\sqrt2\cosh(3ay)}.
\]
The integral
$\int e^{\eta y}/[2\cosh(3ay)]\,dy=(h/2)\sec(h\eta/2)$
and its derivative yield, for $x=h\eta/2$,
\begin{equation}\label{Lsource}
 L_t(\eta):=\int e^{\eta y}\nu_t(y)\,dy
 =\frac3{\sqrt2}\frac{3t\cos x+\sqrt3(2-t)\sin x}
                            {[1+2\cos(2x)]^2},\quad |\Re\eta|<2a.
\end{equation}
The inverse change is
$z=-2\sin x/(\sqrt3\cos x-\sin x)$ and
$r(z)=3/(\sqrt3\cos x-\sin x)^2$.
Substitution gives
$r(z)^{-3/2}L_t(\eta(z))=(t-z)/[\sqrt2(1-z)^2]$.
Since $\sqrt{h_i}=i!d_i/\sqrt2$, coefficient comparison proves
$\ell_t(g)=\int\nu_tg$ on polynomials.

Set
$\rho_t(y)=\int e^{au}\varphi(u)\nu_t(y-u)\,du$.
Exponential tails justify Fubini for each polynomial and show
$\ell_t(T\psi_i)=\int\psi_i\rho_t$.
Its Laplace transform on $-2a<\Re\eta<a$ is
\[
 \int e^{\eta y}\rho_t(y)\,dy
   =\frac{L_t(\eta)}{1+2\cos(h\eta+\pi/3)}.
\]
Applying the same sech integral and derivative at the shifts
$2a,a,0,-a$ shows that its inverse is
\begin{equation}\label{rhosource}
 \rho_t(y)=\frac{A_te^{2ay}+(-y-t+3/2)e^{ay}/\sqrt2+C_t
                +[(1-t)y/2-t+3/2]e^{-ay}/\sqrt2}{2\cosh(3ay)}.
\end{equation}
Fourier uniqueness identifies this expression with the convolution.
The identity
\[w_y=p\cosh(ay)/(2\cosh(3ay))\]
shows that its quotient by $w_y$ is
exactly \eqref{qsource}. At positive infinity
$\rho_t(y)=A_te^{-ay}+O((1+y)e^{-2ay})$; at negative infinity
$\rho_t(y)=O((1+|y|)e^{2ay})$. As
$w_y(y)\sim(y^2/2)e^{-2a|y|}$, one has
$|\rho_t|^2/w_y=O(y^{-2})$ on the positive tail and exponential
decay on the negative one. Thus $q_t\in\mathcal H$.
The basis change and finite inverse formula now prove
\eqref{boundaryexact}. The norm \eqref{Hnorm} also follows by multiplying
the exact recurrence
$H_j/H_{j-1}=3j(j+2)(3j+2)(3j+4)/[4(2j+1)(2j+3)]$, $H_0=2$.

Finally set $A_*=3\sqrt3/4=e^{2\alpha}$. Stirling in \eqref{Hnorm}
gives
$\sqrt{h_j/H_j}\sim [8/(9\,3^{1/4})]A_*^{-j}$.
Equation \eqref{calibrationequiv} gives
$cZ_n(c)/Z_{n+1}(c)\sim[c/(\sqrt2 A_*)]A_*^{-n}$.
Substitution in \eqref{boundaryexact}, with $d_{n-1}\sim n$,
proves \eqref{calcoefficient}, including its alternating sign.
\end{proof}

\subsection{Stable inverses and the analytic coefficient reduction}
For completeness, the finite nonnormal inverses are uniformly bounded
on compact subsets of $\Omega$. In the unweighted representation,
$\widetilde{\mathcal B}=C_0+\mathcal E$, and the earlier proof gives
\[
 \delta(I+\delta P_nC_0P_n)^{-1}P_n\mathcal EP_n
 \longrightarrow\delta(I+\delta C_0)^{-1}\mathcal E
\]
in Hilbert--Schmidt norm, locally uniformly. Nonvanishing of
$\mathcal A$ makes $I+\delta(I+\delta C_0)^{-1}\mathcal E$
invertible. The exact factorization of $I+\delta P_n\widetilde{\mathcal B}P_n$
then bounds its inverse for all sufficiently large $n$; each of the
remaining finitely many sizes is controlled by the negative-root
property. Taking adjoints gives the corresponding bound for
$\mathcal B^*$. The resulting strong convergence on fixed vectors
still does not control a coefficient after division by
$a_{n-1}^{\rm src}=O(1/n)$.

Every $F_n$ is holomorphic and nowhere zero on $\HH_+$, with
$F_n(c)=1$. The earlier interlacing argument gives, for
$R_n(s)=\Log[Z_n(s)/Z_n(c)]$,
\[
 |R_{n+1}(s)-R_n(s)|\le
             \frac{\pi|s-c|}{\sqrt{c\Re s}}.
\]
Thus $F_n$ and $1/F_n$ are locally uniformly bounded. A locally
convergent subsequence and all its derivatives converge by Cauchy's
formula. If every fixed Taylor coefficient at $t=3$ of
$\gamma_{n-1}(t)/a_{n-1}^{\rm src}$ converges, every subsequential
analytic limit has the same Taylor series at $s=c$. Analytic
uniqueness and normal-family compactness then prove full-sequence
locally uniform convergence on $\HH_+$.

For $k\ge1$ the relevant coefficient is exactly
\begin{equation}\label{Taylorcoef}
 \frac{(-1)^k}{a_{n-1}^{\rm src}}\left\{
 \langle\psi_{n-1},(P_n\mathcal B^*P_n)^kP_nq_3\rangle
 -\langle\psi_{n-1},(P_n\mathcal B^*P_n)^{k-1}P_nq^{(1)}\rangle
 \right\}.
\end{equation}
This is the finite Neumann derivative at zero; it involves no
interchange of an infinite series and a size limit. We prove
existence of each fixed-power limit next. The established companion
limit will then evaluate the analytic boundary limit without any
need to compute those powers separately.

\section{Quantitative continuous-Hahn smoothing}\label{sec:smoothing}
The required estimate is a column rate, stronger than compactness.
Let $f_j=\sqrt{w_y}\psi_j\in L^2(\R)$ and $b_y=\sqrt{p(y)}$.
Let $v$ denote the mixed-gamma phase in the exact Jacobi transform:
in the scaled variable $z=y/(3\sqrt3)$,
\[
 G(z)=\Gamma(5/6+iz)\Gamma(7/6-iz),\qquad v=G/|G|.
\]

\begin{theorem}\label{thm:smoothing}
For the kernels in \eqref{kernelreal},
\begin{equation}\label{smoothingrates}
 \sqrt j\|\mathcal E f_j\|_2\longrightarrow0,\qquad
 \sqrt j\|(M_vC_0M_v^*-C_0)f_j\|_2\longrightarrow0.
\end{equation}
The second assertion also holds with input $M_vf_j$ and with the
conjugate phase difference $C_0-M_v^*C_0M_v$. All statements hold
on the full degree sequence, irrespective of the consistent real
leading-sign convention for $\psi_j$.
\end{theorem}

\subsection{An elementary fixed-variable asymptotic, including zero}
We first derive a local asymptotic rather than invoke a large-variable
formula excluding the origin. For positive $A,B$ put $s=A+B$ and
$W(z)=|\Gamma(A+iz)\Gamma(B+iz)|^2$.
The standard continuous-Hahn polynomial $p_j(z;A,B,A,B)$ has positive
leading coefficient and squared norm in $W(z)\,dz$
\begin{equation}\label{Hahnnorm}
 h_j^{A,B}=\frac{2\pi\Gamma(j+2A)\Gamma(j+s)^2\Gamma(j+2B)}
              {(2j+2s-1)\Gamma(j+2s-1)j!}.
\end{equation}
The removable expression at $j=0,s=1/2$ is interpreted by
$(2s-1)\Gamma(2s-1)=\Gamma(2s)$. Neither parameter pair below
requires that convention. The classical generating function
\cite[DLMF 18.23.6]{dlmf} is
\[
 {}_1F_1(A+iz;2A;-it)\,{}_1F_1(B-iz;2B;it)
       =\sum_{j\ge0}\frac{p_j(z)t^j}{(2A)_j(2B)_j}.
\]
Euler's beta integral for ordinary ${}_1F_1$ (DLMF 13.4.1 together
with the normalization conversion 13.2.4) gives exactly
\begin{align}
 p_j(z)&=\frac{i^j\Gamma(j+2A)\Gamma(j+2B)}{j!W(z)}I_j(z),
 \label{Hahnbeta}\\
 I_j(z)&=\int_0^1\!\int_0^1(v-u)^j
 u^{A+iz-1}(1-u)^{A-iz-1}
 v^{B-iz-1}(1-v)^{B+iz-1}\,du\,dv.\nonumber
\end{align}
Absolute convergence and boundedness of the entire exponential on
the integration square justify coefficient extraction. The norm
and polynomial conventions agree with \cite[Equations (1.1)--(1.4)]{cll}.

Uniformly for real $z$ in a compact set, the corners $(u,v)=(0,1)$
and $(1,0)$ yield
\begin{equation}\label{corners}
 I_j(z)=\Gamma(A+iz)\Gamma(B+iz)j^{-s-2iz}
  +(-1)^j\Gamma(A-iz)\Gamma(B-iz)j^{-s+2iz}+O_K(j^{-s-1}).
\end{equation}
To verify the error, at the first corner set $q=1-v$. The smooth
remaining factor equals $1+O_K(u+q)$. On the full triangle
$u,q\ge0$, $u+q\le1$ the leading Dirichlet integral is
\[
 \int u^{A+iz-1}q^{B+iz-1}(1-u-q)^j\,du\,dq
 =\frac{\Gamma(A+iz)\Gamma(B+iz)\Gamma(j+1)}
        {\Gamma(j+s+2iz+1)}.
\]
The $O(u+q)$ term has absolute integral $O_K(j^{-s-1})$.
Removing the portion outside a fixed corner neighborhood costs an
exponentially small error. The second corner is the conjugate version
with $(-1)^j$. Away from these corners $|v-u|\le1-\eps$; the beta
powers have integrable absolute values independent of real $z$.
This proves \eqref{corners}, uniformly through $z=0$.

Let $F_j=\sqrt Wp_j/\sqrt{h_j^{A,B}}$ and
$\eta(z)=\Gamma(A+iz)\Gamma(B+iz)/|\Gamma(A+iz)\Gamma(B+iz)|$.
The gamma-ratio estimate in \eqref{Hahnnorm} gives
\begin{equation}\label{Hahnlocal}
 F_j(z)=\frac1{\sqrt{\pi j}}
 \left\{i^j\eta(z)j^{-2iz}+(-i)^j\overline{\eta(z)}j^{2iz}\right\}
          +O_K(j^{-3/2}).
\end{equation}
This additive uniform error remains valid at zeros and for odd
polynomials at zero. For compactly supported integrable $H$, squaring
and using Riemann--Lebesgue gives
\begin{equation}\label{localmeasure}
 \int H(z)jF_j(z)^2\,dz\longrightarrow\frac2\pi\int H(z)\,dz.
\end{equation}
For compactly supported $L^2$ tests, \eqref{Hahnlocal} similarly
implies $\int H\sqrt jF_j\to0$. The oscillatory frequencies are
$2\log j$ and $4\log j$, respectively.

\subsection{A positive exact norm identity prevents escape of mass}
For the desired pair $(A,B)=(5/6,7/6)$ put $d_0=1/6$ and
$V(z)=|\Gamma(1/6+iz)\Gamma(5/6+iz)|^2$. Gamma recurrence gives
$W(z)=(z^2+d_0^2)V(z)$ exactly. Let $P_j^V,Q_j^W$ be the monic
orthogonal polynomials for $V,W$, with squared norms $H_j^V,H_j^W$.
Both have their degree parity. At $id_0$, the terminating
hypergeometric formula has upper parameter zero, so
\[
 P_j^V(id_0)=\frac{i^j(1/3)_j}{\binom{2j}{j}},\qquad
 r_j=-\frac{P_{j+2}^V(id_0)}{P_j^V(id_0)}
 =\frac{(j+1/3)(j+4/3)(j+1)(j+2)}{4(2j+1)(2j+3)}>0.
\]
The standard leading coefficient in this parameter pair is
$\binom{2j}{j}$. Direct divisibility and orthogonality prove the
quadratic Christoffel relations
\begin{equation}\label{Christoffel}
 Q_j^W(z)=\frac{P_{j+2}^V(z)+r_jP_j^V(z)}{z^2+d_0^2},\qquad
 H_j^W=r_jH_j^V.
\end{equation}
Indeed parity gives the two zeros in the numerator; multiplying by
$z^2+d_0^2$ proves orthogonality against all lower degrees, and the
norm identity follows by pairing with the numerator.

Put $a_j^V=|P_j^V(id_0)|^2/H_j^V$. For $k\le j$ of the same parity,
$P_k^V-[P_k^V(id_0)/P_j^V(id_0)]P_j^V$ is divisible by $z^2+d_0^2$
with quotient of degree less than $j$. Orthogonality and monicity
therefore give
\[
 \langle Q_j^W,P_k^V\rangle_V
       =H_j^V P_k^V(id_0)/P_j^V(id_0).
\]
Opposite-parity pairings vanish. Expand $Q_j^W$ in the $V$ basis
and use \eqref{Christoffel}. For the same normalized function $F_j$
as above this gives the positive exact identity
\begin{equation}\label{normidentity}
 \int\frac{F_j(z)^2}{z^2+1/36}\,dz
       =\frac1{r_ja_j^V}\sum_{\substack{0\le k\le j\\k\equiv j\ (2)}}a_k^V.
\end{equation}
The lowered norm formula \eqref{Hahnnorm} is
$h_j^{1/6,5/6}=2\pi\Gamma(j+1/3)\Gamma(j+5/3)/(2j+1)$,
so
\[
 a_j^V=\frac{(2j+1)\Gamma(j+1/3)}
                  {2\pi\Gamma(1/3)^2\Gamma(j+5/3)}
      \sim\frac{j^{-1/3}}{\pi\Gamma(1/3)^2}.
\]
For either parity, elementary power-sum comparison yields
\[
 \sum_{\substack{k\le j\\k\equiv j\ (2)}}a_k^V
       \sim\frac{3j^{2/3}}{4\pi\Gamma(1/3)^2},\qquad r_j\sim j^2/16.
\]
Thus \eqref{normidentity} proves
\begin{equation}\label{totalmassz}
 j\int\frac{F_j(z)^2}{z^2+1/36}\,dz\longrightarrow12.
\end{equation}
This uses a positive sum and no cancellation of asymptotic terms.

Let $k_0=3\sqrt3$, $y=k_0z$. Up to the real leading sign,
$f_j(y)=k_0^{-1/2}F_j(y/k_0)$ and
$p(y)=k_0^2(z^2+1/36)$. Therefore
\begin{equation}\label{totalmass}
 \|\sqrt j f_j/b_y\|_2^2\longrightarrow4/9.
\end{equation}
On each bounded interval \eqref{localmeasure} gives the limiting
weighted mass $(2/(\pi k_0))\int dy/p(y)$, whose total over the real
line is also $4/9$. Subtracting the compact-interval limit from
\eqref{totalmass} and then enlarging the interval proves
\begin{equation}\label{tightness}
 \lim_{R\to\infty}\limsup_{j\to\infty}
      j\int_{|y|>R}\frac{|f_j(y)|^2}{p(y)}\,dy=0.
\end{equation}
The sequence $u_j=\sqrt j f_j/b_y$ is bounded by \eqref{totalmass}
and weakly null by the compactly supported test conclusion of
\eqref{Hahnlocal}, followed by density. Multiplication by a fixed
bounded unimodular function preserves weak nullity and tightness.

\subsection{A local compactness principle and the original kernel}
If $A$ is bounded and $A\one_{[-R,R]}$ is compact for every $R$,
then a bounded, weakly null, $L^2$-tight sequence $u_j$ satisfies
$\|Au_j\|\to0$. Indeed the compact part tends to zero for fixed
$R$, and the other part is at most $\|A\|\|\one_{|y|>R}u_j\|$.
Taking limsup and then $R\to\infty$ proves the assertion.

Split off the rank-one positive-tail kernel
\[
 R_0(y,x)=\tfrac13G_0(y)r_0(x),\quad
 G_0(y)=g_0(y)e^{2ay},\quad
 r_0(x)=h_0(x)e^{-2ax}\one_{x\ge0}.
\]
Both factors lie in $L^2$. We claim that
$A_0=(\mathcal E-R_0)M_{b_y}$, with the multiplier absorbed in
the input kernel, is bounded and locally compact.
For $x,y<0$ set $r_x=e^{2ax}$ and $d_x=\sqrt{1-r_x+r_x^2}$.
The exact formulas in \eqref{kernelreal} are
$g_0(y)=b_y/[\sqrt2(1+r_y)d_y]$ and
$h_0(x)=\sqrt2d_x/b_x$. The negative-quadrant kernel of $A_0$ is
\begin{align*}
 &\one_{x,y<0}(b_y-b_x)\varphi(x-y)\\
 &\quad+\one_{x,y<0}b_y\varphi(x-y)
                  \left[\frac{d_x}{(1+r_y)d_y}-1\right].
\end{align*}
The first is bounded by $|x-y|\varphi(x-y)$ since $b_y$ is
1-Lipschitz, so Schur's test applies. Restricting its input to a
bounded interval makes it square integrable. The second bracket is
$O(e^{2ax}+e^{2ay})$; exponential off-diagonal decay and these half-line
tails make that whole kernel square integrable.

For $x\ge0$ the kernel is
\[
 g_0(y)\left[\varphi(x-y)-\tfrac13e^{-2a(x-y)}\right]h_0(x)b_x.
\]
Here $h_0(x)b_x=O(e^{2ax})$. The bracket is $O(e^{-4a(x-y)})$
for $x\ge y$ and $O(e^{2a(y-x)})$ for $x<y$.
Use $g_0(y)=O(1+|y|)$ at $-\infty$ and
$O((1+y)e^{-4ay})$ at $+\infty$. The resulting bounds on the three
regions $y<0\le x$, $0\le y\le x$, and $0\le x<y$ are respectively
\[
 C(1+|y|)e^{4ay}e^{-2ax},\qquad
 C(1+y)e^{-2ax},\qquad C(1+y)e^{-2ay}.
\]
Each is square integrable, including the triangular regions.
For $x<0\le y$, the reference and rank-one kernels vanish;
the remaining bound $C(1+y)e^{-6ay}e^{2ax}$ is square integrable.
This proves the claim. The local compactness principle and
\eqref{tightness} now show
\begin{equation}\label{minusRrate}
 \sqrt j\|(\mathcal E-R_0)f_j\|_2=\|A_0u_j\|_2\longrightarrow0.
\end{equation}

The exceptional rank-one term uses the calibrated coefficient,
not a noncalibrated asymptotic. At $t=3$, \eqref{qsource} simplifies to
\begin{equation}\label{q3simple}
 q_3(y)=\frac{\sqrt2[\sqrt3e^{ay}-y-3/2]}{p(y)}.
\end{equation}
Put $F=\sqrt{w_y}q_3$. Then
\begin{equation}\label{subtraction}
 b_y[r_0-\sqrt{2/3}F]\in L^2(\R).
\end{equation}
For $y\ge0$, set $\rho=e^{-2ay}$, $d=\sqrt{1-\rho+\rho^2}$.
Exactly $r_0=\sqrt2d/b_y$ and
$F=[\sqrt3-(y+3/2)e^{-ay}]/(b_yd)$, so the left expression in
\eqref{subtraction} is
$\sqrt2(d-1/d)+\sqrt{2/3}(y+3/2)e^{-ay}/d$, which decays
exponentially. On the negative tail $r_0=0$ and $b_yF$ decays
exponentially up to a polynomial; all factors are bounded near zero.
Writing the function in \eqref{subtraction} as $h_0^{\rm sub}$,
weak nullity and \eqref{calcoefficient} give
\[
 \sqrt j\langle r_0,f_j\rangle
 =\sqrt{2/3}\sqrt j\langle F,f_j\rangle
      +\langle h_0^{\rm sub},u_j\rangle\longrightarrow0.
\]
The rank-one norm therefore tends to zero at the required scale.
Together with \eqref{minusRrate} this proves the first assertion of
\eqref{smoothingrates}.

\subsection{The phase correction has a compact weighted kernel}
The exact gamma phase in \cite[Section 8]{r124} satisfies, on $x,y<0$,
\[
 |v(y)/v(x)-1|\le
     \frac{C|x-y|}{(1+|x|)(1+|y|)}.
\]
Consequently for $D_v=M_vC_0M_v^*-C_0$,
\[
 |D_v(y,x)b_x|\le
       \frac{C\one_{x,y<0}|x-y|\varphi(x-y)}{1+|y|}.
\]
The right side is square integrable, by integrating in the difference
variable and then in $y$. Thus $D_vM_{b_y}$ is Hilbert--Schmidt.
Compactness applied to $u_j\rightharpoonup0$ proves
$\sqrt jD_vf_j\to0$. Apply the identical argument to $M_vu_j$ and
to the conjugate phase difference to obtain all the stated variants.
This completes the proof of Theorem~\ref{thm:smoothing}.

\section{A quantitative Jacobi--Hardy column comparison}\label{sec:columns}
Continue with the angular notation \eqref{wang}. Put $\gamma=2/3$,
$N_j=j+3/2$, $c_0=11\pi/12$, and
\[
 v_j^\pm(\theta)=e^{\pm i(N_j\theta-c_0)}/\sqrt{2\pi},\quad
 v_j=v_j^-+v_j^+,\quad Ue_j=u_j,\quad V_\pm e_j=v_j^\pm,\quad V=V_-+V_+.
\]
These synthesis maps are bounded (the wave maps are restricted circle
Fourier synthesis). Set $X=M_{s_0}HU$, $X_0=M_{s_0}V_-$,
$A=X^*X$, $T_0=X_0^*X_0$.

\begin{theorem}\label{thm:columns}
One has $\sqrt j\|(A-T_0)e_j\|_2\to0$.
After restoring the exact basis phase $i^j$, the Toeplitz model has
real symmetric entries $b_{i-k}$ where
\begin{equation}\label{Toeplitzsymbol}
 b_r=\frac1{2\pi}\int_{-\pi/2}^{\pi/2}w(x+\pi/2)e^{irx}\,dx,
 \qquad |b_r|\le C(1+|r|)^{-5/3}.
\end{equation}
For the original $\psi$-coordinate matrix of $\mathcal B$, denoted
$B$, and this Toeplitz matrix $T_b$, one therefore has
\begin{equation}\label{originalcolumns}
 \|(B-T_b)e_j\|=o((j+1)^{-1/2}).
\end{equation}
\end{theorem}

\subsection{An endpoint Fourier lemma}
Suppose $k(\theta,\phi)$ is locally absolutely continuous in $\phi$,
with envelope $a_0(\theta)$ such that
\[
 |k|\le Ca_0(\theta)d(\phi)^{-1/2},\qquad
 |\partial_\phi k|\le Ca_0(\theta)d(\phi)^{-3/2}.
\]
Suppose that for each fixed interior $\theta$, the two scaled
expressions $d(\phi)^{1/2}k$ and $d(\phi)^{3/2}\partial_\phi k$
tend to zero at both endpoints. Then
\begin{equation}\label{Fourierlittleo}
 \sqrt N\int_0^\pi k(\theta,\phi)e^{\pm iN\phi}\,d\phi\longrightarrow0
\end{equation}
pointwise, bounded in absolute value by $Ca_0(\theta)$.
Split at $1/N$ and $\pi-1/N$. The endpoint pieces are
$O(a_0/\sqrt N)$; one integration by parts on the middle gives
the same order from its boundary and derivative terms. For little-o,
first choose fixed small endpoint neighborhoods where the scaled
bounds are arbitrarily small for this $\theta$, and integrate by
parts on the remaining compact interval, giving $O(1/N)$ there.
This proves \eqref{Fourierlittleo} for real $N\to\infty$.
Dominated convergence gives its $L^2(q(\theta)d\theta)$ version
whenever $\int q a_0^2<\infty$.

\subsection{Two explicit Hardy kernels}
Let $f(\theta)=\log\tan(\theta/2)$, so $f'=1/\sin\theta$.
The difference between $H$ and the compressed ordinary
negative-frequency Hardy operator $C_I$ has, apart from a fixed
factor $-i/(2\pi)$, the kernel
\[
 d_H(\theta,\phi)=
 \frac{\sqrt{f'(\theta)f'(\phi)}}{f(\theta)-f(\phi)}
                         -\frac1{\theta-\phi}.
\]
The canceled diagonal is filled continuously. The explicit bounds are
\begin{equation}\label{Hardyest}
 |d_H|\le\frac C{\sqrt{d(\theta)d(\phi)}},\qquad
 |\partial_\phi d_H|\le\frac C{\sqrt{d(\theta)}d(\phi)^{3/2}}.
\end{equation}
For completeness, in comparable left-endpoint scales set $\phi=\theta r$,
$1/2\le r\le2$. After extracting the scale, the canceled quotient
is smooth through $r=1$, and its required derivatives are uniformly
bounded; hence the bounds are $C/\theta$ and $C/\theta^2$.
In noncomparable scales use
$|f(\theta)-f(\phi)|\ge c|\log(\theta/\phi)|$, whose logarithm
is bounded away from zero. The ordinary Cauchy terms obey
\eqref{Hardyest} separately there. At opposite endpoints the denominator
is bounded below by
$c[1+|\log d(\theta)|+|\log d(\phi)|]$; mixed compact/endpoint
regions obey the same estimates. Direct differentiation uses
\[
 \partial_\phi\frac{\sqrt{f'(\theta)f'(\phi)}}{f(\theta)-f(\phi)}
 =\frac{\sqrt{f'(\theta)f'(\phi)}}{f(\theta)-f(\phi)}
       \left[-\frac{\cot\phi}{2}
                    +\frac{\csc\phi}{f(\theta)-f(\phi)}\right].
\]
For fixed interior $\theta$, the logarithmic denominator gives the
two endpoint little-o conditions. Since $\int w/d$ and
$\int w^2/d$ are finite, the Fourier lemma proves
\begin{equation}\label{Hardywaves}
 \sqrt j\|s_0(H-C_I)v_j^\pm\|_2\to0,
 \qquad \sqrt j\|w(H-C_I)v_j^\pm\|_2\to0.
\end{equation}

The commutator $[H,w]$ has, up to the same fixed factor, kernel
\[
 k_w(\theta,\phi)=(w(\phi)-w(\theta))
       \frac{\sqrt{f'(\theta)f'(\phi)}}{f(\theta)-f(\phi)}.
\]
It obeys the Fourier lemma with square-integrable envelope
\begin{equation}\label{commutenvelope}
 a_0(\theta)=C\left\{d(\theta)^{\gamma-1/2}
       +\frac1{\sqrt{d(\theta)}(1+|\log d(\theta)|)}\right\}.
\end{equation}
Here $|w^{(r)}|\le C_rd^{\gamma-r}$ for $r=0,1,2$.
At comparable endpoint scales the canceled quotient is
$O(\theta^{\gamma-1})$ and its derivative is
$O(\theta^{\gamma-2})$. For $\phi\ge2\theta$ near the same
endpoint, use
\[
 \sup_{2\theta\le\phi\le\eps}
 \frac{\phi^\gamma+\theta^\gamma}{\log(\phi/\theta)}
 \le C\left\{\theta^\gamma+(1+|\log\theta|)^{-1}\right\}.
\]
This follows by putting $v=\log(\phi/\theta)$; the term
$\theta^\gamma e^{\gamma v}/v$ has no interior maximum.
The case $\phi\le\theta/2$ follows directly from
$|w(\phi)-w(\theta)|\le C\theta^\gamma$. Opposite endpoints use
the logarithmic sum denominator from \eqref{Hardyest}; mixed
regions give the second summand in \eqref{commutenvelope}.
Differentiation adds $d(\phi)^{-1}$, using the displayed derivative
formula and the derivative bounds for $w$. For fixed $\theta$ the
scaled expressions still tend to zero, so
\begin{equation}\label{commuterate}
 \sqrt j\|[H,w]v_j^-\|_2\to0.
\end{equation}
The borderline envelope square $d^{-1}/(1+|\log d|)^2$ is integrable.

\subsection{The direct wave defect}
The normalized Jacobi expansion proved from the uniform Bessel
formula in \cite[Section 9]{r124} implies, for $R_j=u_j-v_j$,
\begin{equation}\label{Rbounds}
 |R_j(\theta)|\le C\min\{1,[(j+1)d(\theta)]^{-1}\},\quad
 \|R_j\|_2=O(j^{-1/2}),\quad
 \|s_0R_j\|_2=O(j^{-(1+\gamma)/2}).
\end{equation}
On every compact interior interval $\sqrt jR_j$ tends uniformly
to zero; its global $L^2$ norm is bounded. It is therefore weakly
null. The earlier proof shows $[s_0,H]$ compact, so
\[
 \sqrt j\|s_0HR_j\|_2\to0
\]
by $s_0HR_j=H(s_0R_j)+[s_0,H]R_j$.
This uses compactness on a bounded weakly null sequence with the
scale already included.
The ordinary compressed Hardy tail estimate is
\[
 |C_Iv_j-v_j^-|\le C\{L(N_j\theta)+L(N_j(\pi-\theta))\},\qquad
 L(r)=\begin{cases}1+|\log r|,&r\le1,\\1/r,&r\ge1.\end{cases}
\]
Integration against $w=O(d^\gamma)$ makes its weighted norm
$O(j^{-(1+\gamma)/2})$. Together with \eqref{Hardywaves} and
\eqref{Rbounds}, this proves
\begin{equation}\label{directdefect}
 \sqrt j\|(X-X_0)e_j\|_2\to0.
\end{equation}

\subsection{The adjoint defect must be controlled separately}
We prove
\begin{equation}\label{dualdefect}
 \sqrt j\|(U-V)^*(wv_j^-)\|_{\ell^2}\to0.
\end{equation}
First, uniformly in $k$,
\begin{equation}\label{variation}
 \Var(wR_k)\le C,\qquad (wR_k)(0)=(wR_k)(\pi)=0.
\end{equation}
Here no uncontrolled remainder is differentiated. The exact angular
Jacobi derivative identity for parameters $\alpha_J,\beta_J$ is
\[
 u_k'=q_J(\theta)u_k-
       \sqrt{k(k+\alpha_J+\beta_J+1)}\,
                      u_{k-1}^{(\alpha_J+1,\beta_J+1)},
\]
where $q_J=[(2\alpha_J+1)\cot(\theta/2)
-(2\beta_J+1)\tan(\theta/2)]/4$.
Apply the separate uniform first-order expansions to the two families.
The shifted family has the same $N_k$ and phase $\Phi_k-\pi/2$;
the leading derivative terms cancel, while
$\sqrt{k(k+3)}-N_k=O(1/N_k)$. Consequently
\[
 |R_k'|\le C[d^{-1}+(N_kd^2)^{-1}]\quad(N_kd\ge1),\qquad
 |R_k'|\le CN_k\quad(N_kd<1).
\]
The latter uses the exact identity and endpoint powers $11/6,7/6$,
both greater than one. Integrating $w'R_k+wR_k'$ with
\eqref{Rbounds} proves \eqref{variation}.

The first Jacobi correction is
\[
 A_J(\theta)=\frac{(1/4-\alpha_J^2)\cot(\theta/2)
                  -(1/4-\beta_J^2)\tan(\theta/2)}4,
 \quad (\alpha_J,\beta_J)=(4/3,2/3),\quad\Phi_k=N_k\theta-c_0.
\]
The normalized expansion, with no $1/k$ amplitude term, gives
\begin{equation}\label{L1first}
 \|w[kR_k-\sqrt{2/\pi}A_J\sin\Phi_k]\|_1\longrightarrow0.
\end{equation}
Indeed $wA_J\in L^1$. On $d<1/k$ both terms have integral
$O(k^{-\gamma})$; on the complement the uniform
$O((N_kd)^{-2})$ remainder, multiplied by $kw$, has the same bound.
Replacing $k/N_k$ by one adds $O(k^{-1})$.

Set $\beta_{kj}=\langle R_k,wv_j^-\rangle$. Integration by parts
in \eqref{variation} gives $|\beta_{kj}|\le C/(j+1)$ uniformly
in $k$. By \eqref{L1first}, $k\beta_{kj}$ is, uniformly in $j$,
a fixed linear combination of Fourier coefficients of $wA_J$ at
frequencies $k-j$ and $-(k+j+3)$, plus an error tending to zero
as $k\to\infty$. Fix $\eps>0$ and $M$. The contribution to
$j\sum_k|\beta_{kj}|^2$ from $k<\eps j$ is at most $C\eps$.
The window $|k-j|\le M$ contributes $O(M/j)$.
On $k\ge\eps j$ outside that window, Riemann--Lebesgue and the
uniform error in \eqref{L1first}, together with
$\sum_{k\ge\eps j}k^{-2}=O((\eps j)^{-1})$, give a limsup
bounded by $C\eps^{-1}\sup_{|r|>M}|\widehat{wA_J}(r)|^2$.
Take $j\to\infty$, then $M\to\infty$, then $\eps\downarrow0$.
This proves \eqref{dualdefect}.

Similarly $w$ vanishes at both endpoints and $w'\in L^1$.
Integration by parts and Riemann--Lebesgue give Fourier coefficients
$o(1/|r|)$ uniformly in the tail. The entries of $V_+^*(wv_j^-)$
have frequencies $k+j+3$, so summing their squares gives
\begin{equation}\label{plusdefect}
 \sqrt j\|V_+^*(wv_j^-)\|_{\ell^2}\to0.
\end{equation}
The exact decomposition
\[
 (H-I)wv_j^-=[H,w]v_j^-+w(H-C_I)v_j^-+w(C_I-I)v_j^-
\]
and \eqref{Hardywaves}--\eqref{commuterate} yield
\begin{equation}\label{weightedprojection}
 \sqrt j\|(H-I)(wv_j^-)\|_2\to0.
\end{equation}
The last ordinary Hardy term is $O(1/j)$ in $L^2$ because $2\gamma>1$.

Now put $D=X-X_0$. Exactly
\[
 A-T_0=X^*D+D^*X_0,
\]
\[
 D^*X_0e_j=U^*(H-I)(wv_j^-)
       +(U-V)^*(wv_j^-)+V_+^*(wv_j^-).
\]
Equations \eqref{directdefect}, \eqref{dualdefect},
\eqref{plusdefect}, and \eqref{weightedprojection} prove the first
claim of Theorem~\ref{thm:columns}.

\subsection{Returning to the original operator and restoring phase}
Let $S$ be unitary scaling to $z=y/(3\sqrt3)$, $\mathcal F$ the
unitary Fourier transform with kernel $e^{-iz\xi}$, and $J$ the
Fourier-to-angular coordinate unitary. With the scaled interpretation
of $v$, the exact transform $L=J\mathcal FS M_v$ satisfies
\[
 Lf_j=i^ju_j,\qquad
 J\mathcal FS C_0S^*\mathcal F^*J^*=HM_wH.
\]
The latter identity has no gamma phase. Therefore
\[
 L\widetilde{\mathcal B}f_j-(HM_wH)Lf_j
   =L\mathcal Ef_j+J\mathcal FS(M_vC_0-C_0M_v)f_j.
\]
The second term is the unitary image of
$(M_vC_0M_v^*-C_0)M_vf_j$. The two precise smoothing variants
in Theorem~\ref{thm:smoothing} prove an $o(j^{-1/2})$ error.
Let $Q=\diag(i^j)$. The original coordinate matrix is
$Q^*U^*L\widetilde{\mathcal B}L^*UQ$, and the Hardy model is
$Q^*AQ$. Column norms are unchanged by the diagonal phases.
The raw wave Gram entry is
$(2\pi)^{-1}\int_0^\pi w(\theta)e^{i(i-k)\theta}\,d\theta$.
Conjugating by $Q$ multiplies it by $i^{k-i}$; setting
$x=\theta-\pi/2$ proves exactly \eqref{Toeplitzsymbol} and
\eqref{originalcolumns}.

The corresponding circle symbol is $w(x+\pi/2)$ on
$|x|\le\pi/2$ and zero elsewhere. It is endpoint $C^{2/3}$,
with first and second derivative bounds $d^{-1/3},d^{-4/3}$.
Integrate once by parts, split at distance $1/|r|$ from each
support endpoint, and integrate the interior part once more by
parts. This proves $b_r=O(|r|^{-5/3})$.
The symmetry $w(\pi-\theta)=w(\theta)$ makes $b_r$ real and even.
The absolutely convergent Fourier series vanishes at $x=\pi$, so
\begin{equation}\label{alternatingsum}
 \sum_{r\ge1}(-1)^rb_r=-b_0/2=\sqrt3/4-1/2,
 \qquad b_0=1-\sqrt3/2.
\end{equation}
The last value is the exact scalar mean already evaluated in
\cite[Section 9]{r124}. All parts of Theorem~\ref{thm:columns}
are now proved.

\section{Fixed-source closure and the finite-section edge theorem}\label{sec:edge}
\subsection{An invariant density class}
Choose $R$ with $2h<R<5h/2$. Let $\mathcal C$ consist of functions
$g$ whose density $\rho_g=w_yg$ is holomorphic on $|\Im y|<R$ and,
on every closed narrower strip, satisfies for some integer $M\ge0$
\begin{align}
 \rho_g(y)&=c_ge^{-ay}+O((1+\Re y)^Me^{-2a\Re y})
                                  &&(\Re y\to+\infty),\label{densityclass}\\
 \rho_g(y)&=O((1+|\Re y|)^Me^{2a\Re y})
                                  &&(\Re y\to-\infty).\nonumber
\end{align}
Such functions lie in $\mathcal H$ by the same tail calculation
used for \eqref{qsource}. Cauchy's formula supplies the corresponding
fixed-derivative estimates on narrower strips.

Both $q_3$ and $q^{(1)}$ belong to this class. In
\eqref{rhosource} the affine numerator vanishes identically in $t$
at $y=\pm ih/2,\pm3ih/2$, as direct substitution shows. These
remove the first four zeros of $2\cosh(3ay)$; the next possible
poles are $\pm5ih/2$. The tails give
$c_{q_3}=\sqrt{3/2}$ and $c_{q^{(1)}}=c_{q_3}/4$.

The exact adjoint density formula is
\begin{equation}\label{adjointdensity}
 \rho_{\mathcal B^*g}(y)
   =e^{ay}\int_\R\varphi(y-x)e^{-ax}m(x)\rho_g(x)\,dx.
\end{equation}
It follows directly from the unweighted kernel, first on the real
axis by absolute convergence. Write $F(x)=e^{-ax}m(x)\rho_g(x)$.
It is holomorphic on $|\Im x|<3h/2$, the nearest poles of $m$,
and its horizontal-strip tails are
\[
 F(x)=O(\operatorname{poly}(x)e^{-4a\Re x})\quad(+\infty),\qquad
 F(x)=O(\operatorname{poly}(|x|)e^{a\Re x})\quad(-\infty).
\]
The kernel is holomorphic for $|\Im u|<h$ and has tails
$O(e^{-2a|\Re u|})$ on narrower strips. To extend
\eqref{adjointdensity} to $|\Im y|<5h/2$, choose a contour height
$v$ with $|v|<3h/2$ and $|\Im y-v|<h$. Overlaps agree by contour
shifting; the vertical sides vanish by the displayed tails.
A finite cover gives uniform estimates on each narrower closed strip.

The positive leading coefficient is exactly
\begin{equation}\label{leadingmoment}
 c_{\mathcal B^*g}=\tfrac13\int_\R e^{ax}m(x)\rho_g(x)\,dx.
\end{equation}
To see the error rigorously, subtract $e^{-2a(y-x)}/3$ from the
kernel over the entire chosen contour. For $\Re x\le\Re y$ its
remainder is $O(e^{-4a(\Re y-\Re x)})$; for $\Re x>\Re y$
bound the kernel and subtracted exponential separately. Combining
with the tails of $F$ gives a polynomial times $e^{-4a\Re y}$
before multiplication by $e^{ay}$. The borderline integration over
$0<\Re x<\Re y$ only adds one polynomial power. The final remainder
is $O(\operatorname{poly}(y)e^{-3a\Re y})$, stronger than needed.
The moment integrand in \eqref{leadingmoment} is analytic in the
contour strip and decays at both ends, so its value is contour
independent. On the negative tail split at $\Re x=\Re y$ and
$\Re x=0$. The convolution is
$O(\operatorname{poly}(|y|)e^{a\Re y})$ before multiplication,
giving the required $e^{2a\Re y}$ bound. Thus
\begin{equation}\label{classclosure}
 \mathcal B^*\mathcal C\subset\mathcal C.
\end{equation}

Every member has the same leading coefficient shape:
\begin{equation}\label{shape}
 \langle\psi_j,g\rangle
   =\frac{c_g}{c_{q_3}}a_j^{\rm src}+o(1/(j+1)).
\end{equation}
Here is a proof of the little-o claim, including the endpoint domain.
Subtract $(c_g/c_{q_3})q_3$ and call the result $d_g$. Its density
has two-sided $\operatorname{poly}(y)e^{-2a|\Re y|}$ tails in the
strip. In $z=y/(3h)$ put
$G^\sharp(z)=\Gamma(5/6-iz)\Gamma(7/6+iz)$.
The exactly normalized Fourier input in the Jacobi transform is
$\sqrt{c_W}G(z)d_g(3hz)$, $c_W=81\sqrt3/(8\pi^2)$.
Since $w_y(3hz)$ is a constant times $G(z)G^\sharp(z)$, this input
is a fixed constant times $\rho_{d_g}(3hz)/G^\sharp(z)$.
The reciprocal gamma product is entire. Stirling on horizontal strips
makes this quotient, and all its fixed polynomial multiples, exponentially
decaying. Its holomorphy strip contains $\Im z=\pm2/3$, because
$R>2h$. Contour shifting for its Fourier transform $H_g$ yields
\[
 |H_g(\xi)|+|H_g'(\xi)|\le Ce^{-2|\xi|/3}.
\]
The angular image
$g_\theta(\theta)=\sqrt{2/\sin\theta}\,H_g(\xi(\theta))$
therefore satisfies
$g_\theta=O(d^{5/6})$ and $g_\theta'=O(d^{-1/6})$.
It belongs to $H_0^1(0,\pi)$ and
$\int |g_\theta|^2/d^2<\infty$. Hence it is in the quadratic-form
domain of the Friedrichs Jacobi angular operator, with potential
\[
 \frac{\alpha_J^2-1/4}{4\sin^2(\theta/2)}
 +\frac{\beta_J^2-1/4}{4\cos^2(\theta/2)},
 \qquad (\alpha_J,\beta_J)=(4/3,2/3).
\]
This is the realization with eigenfunctions $u_j$ and eigenvalues
$(j+3/2)^2$, also at the limit-circle endpoint.
Its form expansion gives
$\sum_j(j+3/2)^2|\langle\psi_j,d_g\rangle|^2<\infty$;
the unitary changes coefficients only by phases $i^j$.
Thus $j\langle\psi_j,d_g\rangle\to0$, proving \eqref{shape}.
In particular these coefficients are $O((j+1)^{-1})$ and their
projection-tail norms are $O(n^{-1/2})$.
By \eqref{classclosure}, all fixed iterates of both $q_3$ and
$q^{(1)}$ have this shape. The constants in \eqref{leadingmoment}
retain the full operator, including its compact and rank-one pieces.

\subsection{An abstract edge theorem with nonlocal tails retained}
\begin{theorem}\label{thm:edge}
Let $B$ be bounded on $\ell^2(\N_0)$ and $T_b=T_b^*$ a Toeplitz
operator with entries $b_{i-k}$, where for some $\eps>0$
\[
 |b_r|\le C(1+|r|)^{-1-\eps},\qquad
 \|(B-T_b)e_i\|=o((i+1)^{-1/2}).
\]
Let $q$ be fixed and $u_\ell=(B^*)^\ell q$. Suppose for each fixed
$\ell\ge0$ that
\[
 (u_\ell)_i=\kappa_\ell a_i+o((i+1)^{-1}),\qquad
 a_i\sim c_*( -1)^i/(i+1),\quad c_*\ne0.
\]
Then for every fixed $\ell\ge0$ the full-sequence limit
\[
 \frac{[(P_nB^*P_n)^\ell P_nq]_{n-1}}{a_{n-1}}
\]
exists. Every fixed right-edge profile coordinate has a limit as well.
\end{theorem}
\begin{proof}
Put $Q_n=I-P_n$, $A_n=P_nB^*P_n$, and
$e_{n,\ell}=A_n^\ell P_nq-P_nu_\ell$. Exactly
\begin{equation}\label{defectrecurrence}
 e_{n,0}=0,\qquad
 e_{n,\ell+1}=P_nB^*e_{n,\ell}-P_nB^*Q_nu_\ell.
\end{equation}
The coefficient hypothesis gives $\|Q_nu_\ell\|=O(n^{-1/2})$,
and boundedness of $B$ then gives
$\|e_{n,\ell}\|=O(n^{-1/2})$ for each fixed $\ell$.
Set $K=B-T_b$. If $\|v_n\|=O(n^{-1/2})$, the column hypothesis
implies, uniformly for $i\ge\eta n$, $\eta>0$ fixed,
\begin{equation}\label{compacttail}
 |(K^*v_n)_i|\le\|Ke_i\|\|v_n\|=o(1/n).
\end{equation}
Uniformity follows because
$\sup_{i\ge\eta n}\sqrt{i+1}\|Ke_i\|\to0$.
This uses columns of $K$ to control coordinates of $K^*$, without
assuming an adjoint-column estimate.

Induction in \eqref{defectrecurrence} gives the stronger pointwise bound
\begin{equation}\label{pointwisedefect}
 \sup_{\eta n\le i<n}n|(e_{n,\ell})_i|\le C_{\ell,\eta}.
\end{equation}
For its Toeplitz part split the input at $\eta n/2$. On the upper
part use the preceding inductive bound and $\sum|b_r|<\infty$.
On the lower part the distance is at least $\eta n/2$;
Cauchy--Schwarz bounds the Toeplitz row tail by
$O(n^{-1/2-\eps})$, which times $\|e_{n,\ell}\|=O(n^{-1/2})$
is $O(n^{-1-\eps})$. The source tail $Q_nu_\ell$ has coordinates
$O(1/n)$ for all $i\ge n$, so its Toeplitz image is $O(1/n)$.
The two $K^*$ terms satisfy \eqref{compacttail}. This proves
\eqref{pointwisedefect}, starting from zero.

For each fixed $r\ge0$ let
$E_\ell(r)=\lim_{n\to\infty}(e_{n,\ell})_{n-1-r}/a_{n-1}$
when it exists. It is zero for $\ell=0$. Assuming it exists at
level $\ell$, the exact recurrence gives
\begin{equation}\label{edgeprofile}
 E_{\ell+1}(r)=\sum_{s\ge0}b_{s-r}E_\ell(s)
       -\kappa_\ell\sum_{m\ge0}b_{-r-1-m}(-1)^{m+1}.
\end{equation}
To justify all tails, on input indices $k\ge n/2$ the normalized
defect coordinates are uniformly bounded by \eqref{pointwisedefect};
absolute summability permits a fixed-distance cutoff followed by its
removal. On $k<n/2$, the preceding row-tail estimate is
$O(n^{-1-\eps})=o(a_{n-1})$. On the exterior tail $k=n+m$, the
normalized source coordinates are uniformly bounded and converge,
for each fixed $m$, to $\kappa_\ell(-1)^{m+1}$. Dominated summation
therefore applies to the second series. The $K^*$ terms are
$o(a_{n-1})$ by \eqref{compacttail}. Each limiting profile is bounded
by \eqref{pointwisedefect} with $\eta=1/2$.
This proves existence inductively and gives
\[
 \lim_{n\to\infty}\frac{[A_n^\ell P_nq]_{n-1}}{a_{n-1}}
       =\kappa_\ell+E_\ell(0).
\]
There is no infinite-power interchange in this argument.
\end{proof}

\subsection{Closing the analytic boundary limit}
Apply Theorem~\ref{thm:edge} to the original $B$ in the $\psi$ basis,
using Theorem~\ref{thm:columns} with $\eps=2/3$ and the source shape
\eqref{shape}. Apply it separately to $q_3$ and $q^{(1)}$.
Every coefficient in \eqref{Taylorcoef} therefore converges along
the full sequence. The normal-family argument following that formula
proves $F_n\to F$ locally uniformly on $\HH_+$.
The companion identity gives exactly
\[
 F_n(s)F_{n+1}(s)
  =(s/c)^2\frac{D_n(s^2)/D_n(3)}{D_{n+1}(s^2)/D_{n+1}(3)}.
\]
Since $E_n\to E$ locally uniformly, the right side tends to
$(s/c)^2/P(s)$. Thus $F(s)^2=(s/c)^2/P(s)$.
It is nonzero on $\HH_+$, and $F(c)=1$ selects the branch
$F(s)=(s/c)/\sqrt{P(s)}$. This proves \eqref{Ftarget} without a
residual parity ambiguity.

\subsection{A nonzero leading Galerkin contribution as a check}
There is a useful independent check on the sign and scale.
Define
\[
 \tau_j=\frac{\langle B e_j,Q_{j+1}q_3\rangle}{a_j^{\rm src}},
\]
with source and basis now in coordinate space. The $B-T_b$ part is
$o(1)$ after division by $a_j^{\rm src}$, by
\eqref{originalcolumns} and $\|Q_{j+1}q_3\|=O(j^{-1/2})$.
For the Toeplitz part, the calibration asymptotic makes
$a_{j+r}^{\rm src}/a_j^{\rm src}$ uniformly bounded for $r\ge1$
and gives its fixed-$r$ limit $(-1)^r$. Absolute summability yields
\begin{equation}\label{taulimit}
 \tau_j\longrightarrow\sum_{r\ge1}(-1)^rb_r=\sqrt3/4-1/2.
\end{equation}
Thus the original finite-projection tail is genuinely leading order.
It was retained in \eqref{edgeprofile}; compactness alone would not
have justified discarding it.

\section{Individual amplitudes, inverse thresholds and cumulants}\label{sec:consequences}
\subsection{Proof of the full equivalent}
Write $R_n(s)=\Log[Z_n(s)/Z_n(c)]$ on $\HH_+$, normalized at $c$.
The exact adjacent identity and the now-proved companion limit give
\begin{equation}\label{adjacentsum}
 R_n(s)+R_{n+1}(s)=nf(s^2)+E_n(s^2)+\Log(s/c).
\end{equation}
The exact boundary definition gives
$R_{n+1}(s)-R_n(s)=\Log(s/c)-\Log F_n(s)$.
Because $F_n\to(s/c)e^{-f(s^2)/2}$ locally uniformly and all functions
are nonzero with the same normalization, their analytic logarithms
converge locally uniformly. Hence
\begin{equation}\label{adjacentdifference}
 R_{n+1}(s)-R_n(s)\longrightarrow\tfrac12f(s^2).
\end{equation}
Subtract \eqref{adjacentdifference} from \eqref{adjacentsum} and divide
by two. This yields
\begin{equation}\label{Rlimit}
 R_n(s)-\tfrac n2f(s^2)
 \longrightarrow\tfrac12E(s^2)+\tfrac12\Log(s/c)-\tfrac14f(s^2)
\end{equation}
locally uniformly. Adding the real full-sequence calibration
\eqref{calibrationequiv} proves Theorem~\ref{thm:main}, including its
complex local uniformity. All factors in \eqref{amplitude} are positive
on the positive axis and nonzero on $\HH_+$. There are no zeros or
poles of $P$ there; the definitions at $c$ fix all roots uniquely.
At $s=1$, $P(1)=2/3$, so
$(1/c)^{1/2}P(1)^{-1/4}=2^{-1/4}$. Also
$b+\tfrac12\log(2/3)=\beta$. This verifies \eqref{aequiv}.

\subsection{Amplitude-aware inverse count thresholds}
Let
\[
 C_*=C(1)=C_{\rm cal}2^{-1/4}e^{E(1)/2},\qquad
 N(T)=\min\{n\ge1:a_n\ge T\},\quad y=\log T,\quad r=\sqrt{y/\alpha}.
\]
The minimum exists because \eqref{aequiv} tends to infinity.
The counts are nondecreasing: $A\mapsto\diag(A,1)$ is an injection.
Only large $T$ is considered below.

\begin{corollary}\label{cor:inverse}
Define the real center
\begin{equation}\label{zinverse}
 z_C(T)=r-\frac{\beta}{2\alpha}
       +\frac{\kappa\log r+\beta^2/(4\alpha)-\log C_*}{2\alpha r}.
\end{equation}
There exists a nonnegative function $\eps(T)\to0$ such that
\begin{equation}\label{inversebracket}
 \ceil{z_C(T)-\eps(T)/r}\ \le N(T)\le\ \ceil{z_C(T)+\eps(T)/r}.
\end{equation}
Equivalently the displayed center is
\begin{align}\label{inverseexpanded}
 z_C(T)={}&\sqrt{\frac{\log T}{\alpha}}-\frac{\beta}{2\alpha}
 +\frac{\kappa\log\log T}{4\sqrt{\alpha\log T}}\nonumber\\
 &+\frac{\beta^2/(4\alpha)-\log C_*-(\kappa/2)\log\alpha}
             {2\sqrt{\alpha\log T}}.
\end{align}
Thus the amplitude enters at the $1/\sqrt{\log T}$ scale, and the
remaining error is little-o at that scale before ceilings.
\end{corollary}
\begin{proof}
Set $g(x)=\alpha x^2+\beta x-\kappa\log x+\log C_*$ for $x>0$.
The full equivalent says $\log a_n=g(n)+\eta_n$ with $\eta_n\to0$.
On $[r/2,2r]$, $g'(x)=2\alpha x+\beta-\kappa/x\asymp r$.
Substitution of \eqref{zinverse}, writing its last numerator as
$L_r=\kappa\log r+\beta^2/(4\alpha)-\log C_*$, gives
\[
 g(z_C(T))-y=O((\log r)/r)=o(1).
\]
Indeed the quadratic and linear contributions are
$y-\beta^2/(4\alpha)+L_r+O((\log r)^2/r^2)$,
while $-\kappa\log z_C=-\kappa\log r+O(1/r)$.
The leading terms cancel exactly.

The growth theorem first gives $N(T)\in[r/2,2r]$ for all sufficiently
large $T$ (the endpoints have different fixed quadratic multiples
of $y$). On this interval the tail supremum
$\sup_{n\ge r/2}|\eta_n|$ tends to zero.
Choose $\eps(T)>0$, tending to zero, larger by a fixed sufficient
factor than this supremum plus $|g(z_C)-y|$ and $1/r$.
The lower derivative bound then shows that every integer
$n<z_C-\eps/r$ has $\log a_n<y$, while the integer
$n=\ceil{z_C+\eps/r}$ has $\log a_n\ge y$.
Integers outside $[r/2,2r]$ are handled by the preliminary growth
comparison and monotonicity. These two statements are exactly
\eqref{inversebracket}. Substituting $r=\sqrt{y/\alpha}$ and
$\log r=(\log y-\log\alpha)/2$ proves \eqref{inverseexpanded}.
\end{proof}

The ceilings are essential: even a vanishing real error can decide
which adjacent integer is the threshold when the center is near an
integer. The theorem does not give a fixed explicit decay rate for
$\eps(T)$, nor does it justify discarding it inside a ceiling.

\subsection{Constant terms of every fixed diagonal cumulant}
For fixed $s>0$ assign probability $s^{S(A)}/Z_n(s)$ to a DSASM,
and let $S_n$ denote its number of nonzero diagonal entries.
Put $\mathcal D=s\,d/ds$ and $L(s)=\Log C(s)$, real on the
positive axis. Define
\[
 m(s)=\tfrac12-\frac1{(s+1)(s+2)},\qquad
 v(s)=\frac{s(2s+3)}{(s+1)^2(s+2)^2}>0.
\]
Then $\mathcal DB=m$ and $\mathcal Dm=v$.

\begin{corollary}\label{cor:cumulants}
For each fixed $k\ge1$, locally uniformly on positive fugacity intervals,
\begin{equation}\label{cumulantconstant}
 \operatorname{cum}_k(S_n)=n\mathcal D^{k-1}m(s)
                              +\mathcal D^kL(s)+o(1).
\end{equation}
In particular
\begin{align*}
 \mathbb E_sS_n&=nm(s)+s^2E'(s^2)+\tfrac12-\tfrac12m(s)+o(1),\\
 \Var_s(S_n)&=nv(s)+2s^2E'(s^2)+2s^4E''(s^2)-\tfrac12v(s)+o(1).
\end{align*}
At $s=1$ the linear coefficients are $m(1)=1/3$ and $v(1)=5/36$.
\end{corollary}
\begin{proof}
The locally holomorphic limit \eqref{Rlimit} implies, on a small
complex disk about $z=0$,
\[
 \log\mathbb E_se^{zS_n}
 =n[B(se^z)-B(s)]+L(se^z)-L(s)+o(1)
\]
locally uniformly. Cauchy's formula makes every fixed derivative
of the error $o(1)$. Differentiating proves
\eqref{cumulantconstant}. Inserting
$L(s)=\log C_{\rm cal}+E(s^2)/2+\tfrac12\Log(s/c)
-\tfrac14\Log P(s)$ gives the displayed first two constants.
\end{proof}
This strengthens the bounded cumulant errors in the earlier article.
It relies on complex-local convergence, rather than differentiating
pointwise real asymptotics.

\section{Scope and reproducibility}\label{sec:scope}
The result is a full leading equivalent at each fixed positive
fugacity, with a positive constant characterized by convergent
operators and a Fourier series. Convergence is locally uniform on
the right half-plane. It is not asserted uniformly as $s\to0$ or
$s\to\infty$, nor uniformly in an unbounded derivative order.
No numerical digits for $C(1)$, explicit rate for the little-o terms,
$n^{-2}$ correction, or all-orders conjecture are established here.

The source bundle contains this article's modular TeX, its PDF,
the unchanged earlier-input article's PDF and TeX, a closed SHA-256
inventory, deterministic clean-build and repack scripts, and exact
finite checks with an adversarial input campaign. The finite checks
corroborate identities and normalization; they are not proofs of any
infinite-dimensional estimate, asymptotic convergence or numerical
amplitude. The analytical arguments above establish those limits.
The replay checks package immutability, normal and optimized Python
behavior, independent exact data, two byte-identical clean PDF builds,
and a closed deterministic ZIP round trip. Rebuilding needs only the
locally installed TeX distribution and Python; the scripts make no
network requests and install no dependencies.

\begin{thebibliography}{9}
\bibitem{r124}
\emph{Report 124: Logarithmic asymptotics of diagonally symmetric
alternating sign matrices. Bounded determinant remainders and inverse
thresholds}, 2 October 2026, 34 pp.
Unchanged PDF and complete TeX source included as the earlier-input
appendix to this bundle. Sections 2--4 and 6--14 are the precise
inherited proof dependencies listed in Section~\ref{sec:main}.

\bibitem{bfk}
R. E. Behrend, I. Fischer and C. Koutschan,
\emph{Diagonally symmetric alternating sign matrices},
arXiv:2309.08446v3, 1 October 2026.
\href{https://arxiv.org/abs/2309.08446v3}{arXiv:2309.08446v3}.
The Pfaffian and adjacent determinant are Equations (4.1), (4.11);
calibration product asymptotics are Equations (11.3)--(11.6).

\bibitem{bd}
J. Breuer and M. Duits,
\emph{Central limit theorems for biorthogonal ensembles and asymptotics
of recurrence coefficients}, J. Amer. Math. Soc. 30 (2017), 27--66.
\href{https://doi.org/10.1090/jams/854}{doi:10.1090/jams/854}.
Polynomial right-limit theorem: Corollary 2.2, printed p. 7 of
\href{https://arxiv.org/pdf/1309.6224v2}{arXiv:1309.6224v2}.

\bibitem{cll}
L.-H. Cao, Y.-T. Li and Y. Lin,
\emph{Asymptotic approximations of the continuous Hahn polynomials
and their zeros}, arXiv:1906.03521.
\href{https://arxiv.org/abs/1906.03521}{arXiv:1906.03521}.
Only the polynomial and norm conventions in Equations (1.1)--(1.4)
are used; the local asymptotic needed here is derived in
Section~\ref{sec:smoothing}.

\bibitem{dlmf}
F. W. J. Olver et al., eds., \emph{NIST Digital Library of Mathematical
Functions}, \href{https://dlmf.nist.gov/}{dlmf.nist.gov}.
Continuous-Hahn generating function 18.23.6; ordinary confluent
hypergeometric beta integral via 13.4.1 and 13.2.4; gamma-ratio
asymptotics 5.11.13. Consulted 2 October 2026.

\bibitem{fw}
C. L. Frenzen and R. Wong,
\emph{A uniform asymptotic expansion of the Jacobi polynomials with
error bounds}, Canad. J. Math. 37 (1985), 979--1007.
\href{https://doi.org/10.4153/CJM-1985-053-5}{doi:10.4153/CJM-1985-053-5}.
The exact normalization, first correction, and endpoint/derivative
consequences used here are proved in the supplied Report 124,
Section 9, using DLMF 18.15.6--18.15.9.
\end{thebibliography}

\end{document}
